\documentclass[11pt]{article}
\usepackage[letterpaper,margin=1in]{geometry}
\usepackage{booktabs}
\usepackage{tabularx}
\usepackage{threeparttable}
\usepackage{longtable}
\usepackage{rotating}
\usepackage{caption}
\captionsetup{font=small,labelfont=bf}

\title{Regional Workforce Trends in Texas:\\ A Mid-Decade Assessment}
\author{Office of Economic Opportunity Research}
\date{July 2026}

\begin{document}
\maketitle

\section{Introduction}

Texas continues to add residents and jobs at a pace that outstrips most other
large states, but the gains are not evenly distributed across its regions. This
report reviews recent workforce indicators for a set of illustrative
communities, with attention to labor force participation, occupational demand,
childcare access as a work support, and commuting patterns. All figures in the
tables that follow are simulated for demonstration purposes and should not be
cited as official statistics.

The analysis proceeds in four parts. Section~\ref{sec:labor} summarizes
headline labor market indicators (Table~\ref{tab:headline}) and presents a
detailed community-level panel (Table~\ref{tab:panel}). Section~\ref{sec:occ}
reviews occupational demand projections (Table~\ref{tab:occ}). Section~\ref{sec:supports}
examines childcare capacity as a barrier to parental employment
(Table~\ref{tab:childcare}), including a brief inline comparison of survey
waves. Section~\ref{sec:commute} closes with a wide-format commuting matrix
(Table~\ref{tab:commute}).

\section{Headline labor market indicators}\label{sec:labor}

Statewide, the labor market remained tight through the first half of the year.
Employers in the trade, transportation, and utilities supersector reported the
largest year-over-year gains, while information-sector employment was roughly
flat. Table~\ref{tab:headline} presents the headline indicators for the four
demonstration regions used throughout this report. Two patterns stand out.
First, the gap between the strongest and weakest regional unemployment rates
narrowed relative to the prior year. Second, average weekly earnings growth
was strongest in the region with the youngest workforce, consistent with
compositional shifts toward higher-paying entry occupations.

\begin{table}[htbp]
\centering
\caption{Headline labor market indicators by demonstration region}
\label{tab:headline}
\begin{tabular}{@{}lrrrr@{}}
\toprule
Region & Labor force & Unemp.\ rate (\%) & Avg.\ weekly earnings & YoY jobs (\%) \\
\midrule
Palo Duro Rim    & 412{,}807 & 3.7 & \$1{,}147 & 2.9 \\
Blackland Spine  & 987{,}214 & 4.1 & \$1{,}362 & 3.4 \\
Pineywoods Edge  & 226{,}553 & 5.3 & \$941     & 1.2 \\
Caprock Shelf    & 148{,}039 & 3.1 & \$1{,}018 & 4.6 \\
\bottomrule
\end{tabular}
\end{table}

A fuller community-level panel appears in Table~\ref{tab:panel}, which reports
labor force size, net jobs added, the unemployment rate, and the median hourly
wage for ninety demonstration communities. The panel is deliberately long: it
is intended to exercise page-spanning table handling. Readers interested only
in the regional story can skip directly to Section~\ref{sec:occ} without loss
of continuity.

{\small
\begin{longtable}{@{}lrrrr@{}}
\caption{Community-level workforce panel, demonstration data}
\label{tab:panel} \\
\toprule
Community & Labor force & Net jobs added & Unemp.\ rate (\%) & Median wage \\
\midrule
\endfirsthead
\caption[]{Community-level workforce panel (continued)} \\
\toprule
Community & Labor force & Net jobs added & Unemp.\ rate (\%) & Median wage \\
\midrule
\endhead
\midrule
\multicolumn{5}{r}{\emph{Continued on next page}} \\
\endfoot
\bottomrule
\endlastfoot
Nacogdoches & 3,041 & 113 & 2.1 & \$14.25 \\
Muleshoe & 3,178 & 202 & 2.8 & \$15.55 \\
Dime Box & 3,315 & 291 & 3.5 & \$16.85 \\
Cut and Shoot & 3,452 & 380 & 4.2 & \$18.15 \\
Wink & 3,589 & 469 & 4.9 & \$19.45 \\
Telephone & 3,726 & 558 & 5.6 & \$20.75 \\
Gun Barrel City & 3,863 & 647 & 6.3 & \$22.05 \\
Ding Dong & 4,000 & 736 & 7.0 & \$23.35 \\
Oatmeal & 4,137 & 125 & 2.4 & \$24.65 \\
Noodle & 4,274 & 214 & 3.1 & \$25.95 \\
Bug Tussle & 4,411 & 303 & 3.8 & \$27.25 \\
Uncertain & 4,548 & 392 & 4.5 & \$28.55 \\
Loco & 4,685 & 481 & 5.2 & \$29.85 \\
Notrees & 4,822 & 570 & 5.9 & \$31.15 \\
Happy & 4,959 & 659 & 6.6 & \$32.45 \\
Earth & 5,096 & 748 & 7.3 & \$14.75 \\
Turkey & 5,233 & 137 & 2.7 & \$16.05 \\
Goodnight & 5,370 & 226 & 3.4 & \$17.35 \\
Bigfoot & 5,507 & 315 & 4.1 & \$18.65 \\
Zephyr & 5,644 & 404 & 4.8 & \$19.95 \\
Study Butte & 5,781 & 493 & 5.5 & \$21.25 \\
Point Blank & 5,918 & 582 & 6.2 & \$22.55 \\
North Zulch & 6,055 & 671 & 6.9 & \$23.85 \\
Old Glory & 6,192 & 760 & 2.3 & \$25.15 \\
Pep & 6,329 & 149 & 3.0 & \$26.45 \\
Circle Back & 6,466 & 238 & 3.7 & \$27.75 \\
Fairy & 6,603 & 327 & 4.4 & \$29.05 \\
Lazbuddie & 6,740 & 416 & 5.1 & \$30.35 \\
Salty & 6,877 & 505 & 5.8 & \$31.65 \\
Tarzan & 7,014 & 594 & 6.5 & \$32.95 \\
Nacogdoches South & 7,151 & 683 & 7.2 & \$15.25 \\
Muleshoe South & 7,288 & 772 & 2.6 & \$16.55 \\
Dime Box South & 7,425 & 161 & 3.3 & \$17.85 \\
Cut and Shoot South & 7,562 & 250 & 4.0 & \$19.15 \\
Wink South & 7,699 & 339 & 4.7 & \$20.45 \\
Telephone South & 7,836 & 428 & 5.4 & \$21.75 \\
Gun Barrel City South & 7,973 & 517 & 6.1 & \$23.05 \\
Ding Dong South & 8,110 & 606 & 6.8 & \$24.35 \\
Oatmeal South & 8,247 & 695 & 2.2 & \$25.65 \\
Noodle South & 8,384 & 784 & 2.9 & \$26.95 \\
Bug Tussle South & 8,521 & 173 & 3.6 & \$28.25 \\
Uncertain South & 8,658 & 262 & 4.3 & \$29.55 \\
Loco South & 8,795 & 351 & 5.0 & \$30.85 \\
Notrees South & 8,932 & 440 & 5.7 & \$32.15 \\
Happy South & 9,069 & 529 & 6.4 & \$14.45 \\
Earth South & 9,206 & 618 & 7.1 & \$15.75 \\
Turkey South & 9,343 & 707 & 2.5 & \$17.05 \\
Goodnight South & 9,480 & 796 & 3.2 & \$18.35 \\
Bigfoot South & 9,617 & 185 & 3.9 & \$19.65 \\
Zephyr South & 9,754 & 274 & 4.6 & \$20.95 \\
Study Butte South & 9,891 & 363 & 5.3 & \$22.25 \\
Point Blank South & 10,028 & 452 & 6.0 & \$23.55 \\
North Zulch South & 10,165 & 541 & 6.7 & \$24.85 \\
Old Glory South & 10,302 & 630 & 2.1 & \$26.15 \\
Pep South & 10,439 & 719 & 2.8 & \$27.45 \\
Circle Back South & 10,576 & 808 & 3.5 & \$28.75 \\
Fairy South & 10,713 & 197 & 4.2 & \$30.05 \\
Lazbuddie South & 10,850 & 286 & 4.9 & \$31.35 \\
Salty South & 10,987 & 375 & 5.6 & \$32.65 \\
Tarzan South & 11,124 & 464 & 6.3 & \$14.95 \\
Nacogdoches East & 11,261 & 553 & 7.0 & \$16.25 \\
Muleshoe East & 11,398 & 642 & 2.4 & \$17.55 \\
Dime Box East & 11,535 & 731 & 3.1 & \$18.85 \\
Cut and Shoot East & 11,672 & 120 & 3.8 & \$20.15 \\
Wink East & 11,809 & 209 & 4.5 & \$21.45 \\
Telephone East & 11,946 & 298 & 5.2 & \$22.75 \\
Gun Barrel City East & 3,083 & 387 & 5.9 & \$24.05 \\
Ding Dong East & 3,220 & 476 & 6.6 & \$25.35 \\
Oatmeal East & 3,357 & 565 & 7.3 & \$26.65 \\
Noodle East & 3,494 & 654 & 2.7 & \$27.95 \\
Bug Tussle East & 3,631 & 743 & 3.4 & \$29.25 \\
Uncertain East & 3,768 & 132 & 4.1 & \$30.55 \\
Loco East & 3,905 & 221 & 4.8 & \$31.85 \\
Notrees East & 4,042 & 310 & 5.5 & \$33.15 \\
Happy East & 4,179 & 399 & 6.2 & \$15.45 \\
Earth East & 4,316 & 488 & 6.9 & \$16.75 \\
Turkey East & 4,453 & 577 & 2.3 & \$18.05 \\
Goodnight East & 4,590 & 666 & 3.0 & \$19.35 \\
Bigfoot East & 4,727 & 755 & 3.7 & \$20.65 \\
Zephyr East & 4,864 & 144 & 4.4 & \$21.95 \\
Study Butte East & 5,001 & 233 & 5.1 & \$23.25 \\
Point Blank East & 5,138 & 322 & 5.8 & \$24.55 \\
North Zulch East & 5,275 & 411 & 6.5 & \$25.85 \\
Old Glory East & 5,412 & 500 & 7.2 & \$27.15 \\
Pep East & 5,549 & 589 & 2.6 & \$28.45 \\
Circle Back East & 5,686 & 678 & 3.3 & \$29.75 \\
Fairy East & 5,823 & 767 & 4.0 & \$31.05 \\
Lazbuddie East & 5,960 & 156 & 4.7 & \$32.35 \\
Salty East & 6,097 & 245 & 5.4 & \$14.65 \\
Tarzan East & 6,234 & 334 & 6.1 & \$15.95 \\
\end{longtable}
}

\section{Occupational demand}\label{sec:occ}

Projected occupational demand reflects both replacement needs and net new
growth. Health care support roles and skilled trades dominate the top of the
demand rankings in every region, though the mix varies with local industry
structure. Table~\ref{tab:occ} lists selected occupations with projected
openings and typical education requirements; the description column uses the
full available text width, which is why the table is set with a flexible
column layout.

\begin{table}[htbp]
\centering
\caption{Selected occupations: projected annual openings and entry requirements}
\label{tab:occ}
\begin{tabularx}{\textwidth}{@{}lrrX@{}}
\toprule
Occupation & Openings & Growth (\%) & Notes on entry pathway \\
\midrule
Wind turbine technician & 1{,}943 & 44.8 & Certificate programs at Sweetwater and Borger campuses; median completion time 14 months. \\
Licensed vocational nurse & 8{,}671 & 12.6 & Bridge programs allow articulation to RN within 26 months at partner colleges. \\
Diesel mechanic & 4{,}058 & 9.4 & Registered apprenticeship expanded to 17 new sponsor employers in Lamesa and Sonora. \\
Cybersecurity analyst & 2{,}317 & 31.9 & Stackable credential ladder; entry possible with sub-baccalaureate certificate plus lab hours. \\
Early childhood teacher & 6{,}482 & 7.3 & CDA credential remains the modal entry point; wage supplement pilot active in three regions. \\
\bottomrule
\end{tabularx}
\end{table}

The occupational picture interacts with training capacity. Where community
college programs are oversubscribed, projected openings can go unfilled even
in high-wage fields. Regional workforce boards have responded by funding
evening cohorts and equipment purchases, though instructor recruitment remains
the binding constraint in most rural service areas.

\section{Work supports: childcare capacity}\label{sec:supports}

Access to reliable childcare is one of the most consistent predictors of
sustained labor force participation among parents of young children. Licensed
capacity has grown modestly, but subsidized seats remain scarce relative to
estimated need, especially for infant care. Table~\ref{tab:childcare} reports
demonstration estimates of licensed capacity, estimated need, and the
resulting gap, with methodological caveats in the table notes.

\begin{table}[htbp]
\centering
\begin{threeparttable}
\caption{Licensed childcare capacity and estimated need, children under 6}
\label{tab:childcare}
\begin{tabular}{@{}lrrrr@{}}
\toprule
Region & Licensed seats & Est.\ need & Gap & Gap ratio \\
\midrule
Palo Duro Rim   & 18{,}412 & 27{,}906 & 9{,}494  & 1.52 \\
Blackland Spine & 61{,}377 & 74{,}218 & 12{,}841 & 1.21 \\
Pineywoods Edge & 9{,}063  & 16{,}540 & 7{,}477  & 1.83 \\
Caprock Shelf   & 5{,}891  & 8{,}774  & 2{,}883  & 1.49 \\
\bottomrule
\end{tabular}
\begin{tablenotes}\footnotesize
\item[a] Estimated need assumes 68.5 percent of children under 6 have all
resident parents in the labor force.
\item[b] Gap ratio is estimated need divided by licensed seats; values above
1.0 indicate unmet need. Family, friend, and neighbor care is excluded.
\end{tablenotes}
\end{threeparttable}
\end{table}

Survey evidence points in the same direction. Among parents who left a job in
the past year, the share citing care disruptions as the primary reason has
risen in each of the last three waves, as the following inline tabulation of
wave-over-wave responses shows:
\begin{center}
\begin{tabular}{@{}lccc@{}}
Reason cited        & Wave A & Wave B & Wave C \\
Care disruption     & 23.4\% & 28.9\% & 33.1\% \\
Schedule volatility & 19.7\% & 18.2\% & 17.6\% \\
Relocation          & 11.8\% & 10.4\% & 9.9\%  \\
\end{tabular}
\end{center}
and the increase is concentrated among parents of infants. The pattern holds
after adjusting for industry mix and shift type, suggesting the constraint is
structural rather than cyclical.

\section{Commuting patterns}\label{sec:commute}

Finally, commuting flows knit the demonstration regions together. Because the
flow matrix requires more horizontal space than a portrait page allows, it is
presented in landscape orientation in Table~\ref{tab:commute}. The dominant
flows run from bedroom communities toward the two employment cores, but
reverse commuting has grown as distribution centers locate along interstate
corridors outside the cores.

\begin{sidewaystable}
\centering
\caption{Estimated daily worker flows between demonstration regions (thousands)}
\label{tab:commute}
\begin{tabular}{@{}lrrrrrr@{}}
\toprule
 & \multicolumn{6}{c}{Workplace region} \\
\cmidrule(l){2-7}
Residence region & Palo Duro Rim & Blackland Spine & Pineywoods Edge & Caprock Shelf & Out-of-state & Remote \\
\midrule
Palo Duro Rim   & 171.4 & 38.2  & 4.7   & 12.9 & 3.3  & 21.8 \\
Blackland Spine & 29.6  & 402.1 & 17.3  & 6.4  & 11.2 & 58.7 \\
Pineywoods Edge & 3.1   & 22.8  & 88.6  & 1.9  & 7.4  & 13.2 \\
Caprock Shelf   & 14.7  & 8.3   & 1.2   & 61.9 & 2.6  & 8.1  \\
\bottomrule
\end{tabular}
\end{sidewaystable}

\section{Conclusion}

The demonstration data reviewed here illustrate a consistent story: regional
labor markets are tight, occupational demand is concentrated in middle-skill
fields with constrained training pipelines, and childcare scarcity continues
to suppress parental labor supply. Policy attention to training capacity and
care infrastructure would address the two most binding constraints
simultaneously. Subsequent reports will extend the panel in
Table~\ref{tab:panel} and refine the flow estimates in
Table~\ref{tab:commute} as new demonstration waves become available.

\end{document}
